\justifying
\NSFCBodyText

\subsubsection{与本项目相关的研究工作积累}

申请人颜丰为新疆大学副教授、博士、硕士研究生导师，已开展多尺度时间序列、能源预测和多源信息融合研究。与本项目最直接相关的基础，是从多时间尺度中提取可复用表征、处理不同信息源之间的关系，以及组织模型训练和比较实验。这些能力可用于本课题的资料处理、联合表征和时间连续性分析。

申请人在 DWT-Former 工作中担任通讯作者，研究小波多尺度特征与 Transformer 时间表示的结合，成果发表于 Energy\cite{DWT2025}。该工作支持本课题对缓慢变化背景和局地波动进行分解的技术准备。申请人在 WaveKAN 工作中担任共同第一作者，研究自适应小波分解与长时序建模，成果发表于 Neurocomputing\cite{WaveKAN2026}。该工作为设计不同时间尺度的对照及分析长期变化提供方法积累。

上述成果的研究对象与 CIESM 气候场有所不同，其作用是支撑模型组织与实验方法。现有资料未提供本课题专属的 CIESM 数据产品、训练模型、消融实验或性能日志，因此本申请不列示已经获得的 CIESM 重建指标。第一季度的数据审计与基线实验将承担从既有时序研究向本任务转化的检验。

\subsubsection{已取得的研究工作成绩}

两项代表性成果分别体现申请人在多尺度时间表征和自适应序列建模方面的研究积累，为本课题的时空信息组织、模型比较和训练稳定性分析提供了方法储备。其完整题录见参考文献，既有成果与本课题拟开展的气候实验分别列示。

申请人主持在研的“融合多源数据大模型驱动的光伏发电功率预测方法研究”，并参与“数据驱动下基于互补型生成对抗网络与神经统计模型的公路沙害智能预警研究”，可为多源数据整理和区域相关资料处理提供经验。具体项目事实及与本课题的区别见正在承担项目情况。已有项目任务、资金和成果归属按各自任务书执行，本课题以气候模式场增强为独立研究对象。

\subsubsection{可行性分析}

\textbf{方法基础。}
本项目从统计校正和等信息独立模型起步，逐步检验共享表征、地形条件与约束。既有时序建模经验可支持实现这些模型，公开降尺度研究可提供比较依据。科学问题通过增量消融与独立验证回答，无需在一年内同时从头开发大型气候模式或全球基础模型。

\textbf{资料与算力。}
公开 CIESM 和再分析资料提供先行研究路径，申请人已确认可稳定使用的 GPU 及软件环境支持区域分块训练。公开资料的变量频率和共同覆盖期仍需实际审计，因此资料获取与模型研究分阶段推进，避免在样本尚不成立时预设性能结果。

\textbf{组织安排。}
申请人负责总体设计、模型与验证方法、进度管理和成果组织。本课题的主要参与者及具体分工为\textbf{【待补 F-CIESM-02：参与者名单与分工】}，落实后分别承担资料处理、模型实验和集成评价。其他项目的学生名单不自动列入本课题。

\subsubsection{研究风险与应对措施}

\textbf{数据与频率风险。}
若第一季度发现 CIESM 关键变量缺失或仅有月频资料，及时寻求实验室提供合适版本，保留公开资料的可复现先行试验。月尺度结果只作为阶段性证据，不替代日极端与日内增强考核。再分析的细网格若主要来自插值，则同步报告插值基线，限制对真实新增细节的解释。

\textbf{技术风险。}
若联合训练使某些变量持续退化，采用等信息独立输出或降低共享深度；若软约束与极端结构冲突，保留数学上成立的条件并减弱经验惩罚。判断依据来自固定验证集，修改后重新进行受控比较。若随机模型无法形成可校准的样本分布，交付确定性核心模型及其误差说明，不保留未经检验的概率承诺。

\textbf{进度与应用风险。}
实验室接口对接从第一季度开始，公开格式测试与接口联调分开推进。若正式接收延迟，保留可复现兼容性结果并如实列示未完成的应用环节，继续按约定联调。论文尽早完成主要设计和初稿，预留审稿修改时间；将论文投稿、接收与发表作为不同状态管理。

\textbf{资源风险。}
以区域小块、轻量模型和少量种子重复控制训练规模，记录 GPU 时间和存储占用。若资料规模或生成模型开销高于预期，优先保障四类变量的日尺度核心、等信息对照和集成验证，将日内扩展顺延为后续工作，不依赖尚未落实的额外算力支撑核心目标。
